\documentclass[11pt]{amsart}
\usepackage{amsmath,amssymb,amsthm}
\usepackage[margin=1in]{geometry}
\usepackage{hyperref}
\hypersetup{colorlinks=true, urlcolor=blue, linkcolor=black, citecolor=black}

% Cyrillic Sha
\usepackage[OT2,T1]{fontenc}
\DeclareSymbolFont{cyrletters}{OT2}{wncyr}{m}{n}
\DeclareMathSymbol{\Sha}{\mathalpha}{cyrletters}{"58}

\newtheorem{defn}{Definition}[section]
\newtheorem{thm}[defn]{Theorem}
\newtheorem{theorem}[defn]{Theorem}
\newtheorem{prop}[defn]{Proposition}
\newtheorem{lemma}[defn]{Lemma}
\newtheorem{cor}[defn]{Corollary}
\newtheorem{conj}[defn]{Conjecture}
\newtheorem{remark}[defn]{Remark}
\newtheorem{example}[defn]{Example}
\newtheorem*{ack}{Acknowledgements}

\newcommand{\Q}{\mathbb Q}
\newcommand{\Z}{\mathbb Z}
\newcommand{\F}{\mathbb F}
\newcommand{\kron}[2]{\left(\frac{#1}{#2}\right)}
\DeclareMathOperator{\rank}{rank}
\DeclareMathOperator{\Gal}{Gal}

\setlength{\emergencystretch}{3em}

\begin{document}

\bibliographystyle{plain}
\title[A rank $\geq 34$ curve over a quadratic field]{An elliptic curve of rank $\geq 34$ and root number $-1$ over a quadratic field}

\author{Dylan T. Costa}
\address{University of Connecticut, Department of Mathematics, Storrs, CT 06269, USA}
\email{dylan.costa@uconn.edu}
\urladdr{https://dylan-costa.github.io/}

\author{\'Alvaro Lozano-Robledo}
\address{University of Connecticut, Department of Mathematics, Storrs, CT 06269, USA}
\email{alvaro.lozano-robledo@uconn.edu}
\urladdr{http://alozano.clas.uconn.edu}

\subjclass[2020]{Primary: 11G05, Secondary: 11G40, 14H52.}

\begin{abstract}
In this article, we find a quadratic twist over $\Q$ of rank  $\geq 2$ (and $\geq 3$ assuming parity) of an elliptic curve $E/\Q$ of rank $\geq 32$, recently found by Elkies, Klagsbrun, and Li. We use a family of quadratic twists of rank $\geq 2$ over $\Q(t)$ constructed by Mestre to search for twists $E^D$ with root number $-1$, computing all root numbers without factoring the (very large) twisting parameter $D$. As a result, we exhibit an explicit integer $D$, with $411$ digits (a square times an integer with $381$ digits), such that $E^D$ has root number $-1$ and two independent rational points. In particular, the rank of $E$ over $\Q(\sqrt{D})$ is at least $34$, unconditionally, and at least $35$ if we assume the Parity Conjecture.
\end{abstract}

\maketitle

\section{Introduction}\label{sec-intro}

Let $F$ be a number field and let $E/F$ be an elliptic curve. The Mordell--Weil theorem states that the group $E(F)$ of $F$-rational points on $E$ is finitely generated and, therefore, we have an isomorphism
$$E(F)\cong E(F)_\text{tors}\oplus \Z^{R_{E/F}},$$
where $E(F)_\text{tors}$ is the finite subgroup of points of finite order, and $R_{E/F}=\rank_\Z E(F)\geq 0$ is the rank of $E$ over $F$. It is not known whether the rank of elliptic curves over $\Q$ can be arbitrarily large. Recently, Noam Elkies, Zev Klagsbrun, and Samuel Li have found an elliptic curve $E/\Q$ with $R_{E/\Q}\geq 32$ (see \cite{icarm} and Section \ref{sec-curve} below), which has the largest known rank over $\Q$ as of the writing of this article.

In this article, we investigate rank growth of $E$ when we base extend to a quadratic field $K=\Q(\sqrt{D})$. Let $E^D$ denote the quadratic twist of $E$ by $D$; then it is well known that
\begin{equation}\label{eq-ranksum}
R_{E/K}=R_{E/\Q}+R_{E^D/\Q}
\end{equation}
(see \cite{kida}, or \cite[Ch.~X]{silverman1}); in particular, the rank of $E$ grows from $\Q$ to $K$ if and only if $E^D$ has positive rank. Moreover, if $W(E^D)=-1$ and $E^D$ has two independent rational points, then the Parity Conjecture (Conjecture \ref{conj-parity}) implies that $R_{E^D/\Q}\geq 3$. Mestre has constructed a family of quadratic twists of $E$ with two independent points over $\Q(t)$ (see \cite{mestre}, and also \cite{rubinsilverberg}, \cite{stewarttop}, \cite{gusictadic}), so our task is reduced to finding a specialization $t\in\Q$ for which the twist has root number $-1$. The main difficulty is that the resulting twisting parameters have hundreds of digits that are hard to factor, so the usual computation of $W(E^D)$ via the conductor of $E^D$ is not available. In this paper, we prove the following result:


\begin{thm}\label{thm-main-intro}
Let $E/\Q$ be the elliptic curve of Section \ref{sec-curve}, with short Weierstrass model $y^2 = x^3+ax+b$ as in Section \ref{sec-mestre}, let
\begin{equation*}
   g(t)=-ab\,(t^2+1)\left(b^2(t^4+t^2+1)^3+a^3t^4(t^2+1)^2\right)\in\Z[t],
\end{equation*}
let $D=g(3)$, and let $K=\Q(\sqrt{D})$. Then 
\begin{enumerate}
    \item $W(E/\Q)=1$,
    \item $W(E^D/\Q)=-1$, and $R_{E^D/\Q}\geq 2$.
\end{enumerate}
Thus, $R_{E/K}\geq 34$. Moreover, if we assume the Parity Conjecture for elliptic curves over $\Q$, or $\Sha(E^D/\Q)[p^\infty]$ is finite for some prime $p$, then $R_{E^D/\Q}\geq 3$ and $R_{E/K}\geq 35$.
\end{thm}

The integer $D$ has $411$ digits. After removing square factors of $D$,  we have $K=\Q(\sqrt{D_0})$ where $D_0=5\cdot 29\cdot 67\cdot 191\cdot 269\cdot 521\cdot 709\cdot 37847\cdot C$ is an integer with $381$ digits (see Example \ref{ex-t3}), and $C$ is a composite number with $362$ digits that we have not factored. The proof of Theorem \ref{thm-main-intro} will avoid factoring $C$.

\subsection{Structure of the article} The article is organized as follows. In Section \ref{sec-curve} we describe the curve of Elkies, Klagsbrun, and Li, and we compute its global root number. In Section \ref{sec-mestre} we recall Mestre's family of quadratic twists of rank $\geq 2$. In Section \ref{sec-rootnumbers} we explain how to compute root numbers of twists without factoring the twisting parameter, and in Section \ref{sec-independence} we explain how to prove the independence of the two points without computing heights. In Section \ref{sec-search} we describe our search for twists of root number $-1$ in Mestre's family. Finally, Theorem \ref{thm-main-intro} is shown in Section \ref{sec-proof}.

\subsection{Acknowledgements:}
The authors thank Noam Elkies, Zev Klagsbrun, and Samuel Li for making their curve publicly available. We would like to thank ICARM (which is supported by NSF Grant DMS 2425401). The computations in this article were done with Magma \cite{magma}, and we used Claude  Opus 5.5 to help write the Magma code to carry out the searches and verify the results. The code and data are available at \url{https://github.com/alozanoroble/rank35}; in particular, the Magma file \texttt{g3-twist-data.m} verifies every claim about the case $t=3$ when it is loaded.

\subsection{Declaration of AI use:} We have used Claude Opus 5.5 for literature review, to help us write Magma code, and to help us search for a twisting parameter $D$ such that $E^D$ has the properties in Theorem \ref{thm-main-intro}. Claude also helped putting together a first draft of this paper, but the authors rewrote it for readability. 

\section{The curve of Elkies, Klagsbrun, and Li}\label{sec-curve}

Let $E/\Q$ be the elliptic curve given by $E: y^2+xy=x^3+a_4x+a_6$, with
\begin{align*}
a_4&=-2862392899139891916550046700194750245224\\
&\qquad 0040593445344054331428726899413,\\
a_6&=\phantom{-}16366380970220279630250861194633604823688757308892533\\
&\qquad 34875693440181143871885452358966522018914326964558817.
\end{align*}
This curve was found by Noam Elkies, Zev Klagsbrun, and Samuel Li by specializing a rank $17$ fibration of the known K3 surface of discriminant $1152$, and it was submitted to the ICARM elliptic curve rank leaderboard as curve \#4239 (see \cite{icarm}), together with $32$ rational points. We have verified with Magma that the height regulator of these points is approximately $7.166\cdot 10^{42}\neq 0$ (with Magma's normalization of the canonical height), and that $E(\Q)_\text{tors}$ is trivial. Hence, $R_{E/\Q}\geq 32$.

The model given above is minimal, and the conductor of $E$ is
$$N_E=2\cdot 3\cdot 5^2\cdot 11^2\cdot 13\cdot 17\cdot 23\cdot 31\cdot 37\cdot 41\cdot 43\cdot 53\cdot 89\cdot p\cdot q,$$
where
\begin{align*}
p&=2220713642109042278879558804334283,\\
q&=3500766112881427717056865531427442780770809\\
&\qquad 0794959852326328757722465802911598967036101\\
&\qquad 6516064255191976556239262775289081591134547.
\end{align*}
The curve $E$ has additive reduction at $5$ (of Kodaira type $I_0^*$) and at $11$ (of type $IV$), and multiplicative reduction at all other bad primes. The discriminant and the conductor agree with the values on the ICARM website.

The root number of $E$, denoted $W(E/\Q)$, is not listed in \cite{icarm}, so we have computed it. The local root numbers $W_\ell(E/\Q)$ are as follows:
\begin{center}
\begin{tabular}{|c|ccccccccccccccc|}
\hline
$\ell$ & $2$ & $3$ & $5$ & $11$ & $13$ & $17$ & $23$ & $31$ & $37$ & $41$ & $43$ & $53$ & $89$ & $p$ & $q$\\ \hline
type & $I_{14}$ & $I_9$ & $I_0^*$ & $IV$ & $I_2$ & $I_2$ & $I_2$ & $I_2$ & $I_3$ & $I_5$ & $I_2$ & $I_2$ & $I_2$ & $I_1$ & $I_1$\\
\hline
$W_\ell(E/\Q)$ & $-1$ & $-1$ & $1$ & $-1$ & $1$ & $1$ & $1$ & $1$ & $-1$ & $-1$ & $1$ & $1$ & $1$ & $1$ & $1$\\
\hline
\end{tabular}
\end{center}
Since $W_\infty=-1$ for every elliptic curve over $\Q$, we compute the global root number $W(E/\Q)$ as
$$W(E)=W_\infty(E/\Q)\prod_{\ell\mid N_E}W_\ell(E/\Q)=(-1)\cdot(-1)^5=1.$$
Thus, the Parity Conjecture below predicts that $R_{E/\Q}$ is even.

\begin{conj}[Parity Conjecture]\label{conj-parity}
Let $E/\Q$ be an elliptic curve. Then $W(E/\Q)=(-1)^{R_{E/\Q}}$.
\end{conj}

\section{Mestre's family of quadratic twists of rank $\geq 2$}\label{sec-mestre}

Since we are working over $\Q$, the curve $E$ can be put into a short Weierstrass model. Explicitly, the change of variables $(x,y)\mapsto(36x+3,\,108(2y+x))$ takes the model of Section \ref{sec-curve} to a short Weierstrass model, which we write again as $E: y^2 = x^3+ax+b$, where $a=-27c_4$ and $b=-54c_6$ are given by
\begin{align*}
a&=-37096611972852999238488605234523963178103092609105165894413531630061639275,\\
b&=76358987054659736642898417989682546776692102018928001307426168848663\\
&\qquad 338020996943087257548997707433748641083750.
\end{align*}
 We define
\begin{equation}\label{eq-g}
g(t)=-ab\,(t^2+1)\left(b^2(t^4+t^2+1)^3+a^3t^4(t^2+1)^2\right)\in\Z[t],
\end{equation}
a polynomial of degree $14$. Mestre showed in \cite{mestre} that the quadratic twist 
$$E_{g}\colon y^2=x^3+ag(t)^2x+bg(t)^3$$ has rank at least $2$ over $\Q(t)$, and we shall use the explicit form of the construction given in \cite[Example 4.1]{gusictadic} (see also \cite[Theorem 3.7]{rubinsilverberg} and \cite[Theorem 3]{stewarttop}), with two independent points
\begin{equation}\label{eq-PQ}
\begin{aligned}
P(t)&=\left(-\frac ba\cdot\frac{(t^2+t+1)(t^2-t+1)}{t^2+1}\cdot g(t),\ \frac{g(t)^2}{a^2(t^2+1)^2}\right),\\
Q(t)&=\left(-\frac ba\cdot\frac{(t^2+t+1)(t^2-t+1)}{t^2(t^2+1)}\cdot g(t),\ \frac{g(t)^2}{a^2t^3(t^2+1)^2}\right).
\end{aligned}
\end{equation}
For $t_0\in\Q$ with $t_0\neq 0,\pm 1$ and $g(t_0)\neq 0$, we write $D(t_0)$ for the integer $\operatorname{num}(g(t_0))\cdot\operatorname{den}(g(t_0))$, which is in the same square class as $g(t_0)$. Equivalently, if $t_0=m/n$ in lowest terms, then $D(t_0)$ is in the same square class as the integer
$$G(m,n)=n^{14}g(m/n)=-ab\,(m^2+n^2)\left(b^2(m^4+m^2n^2+n^4)^3+a^3m^4n^4(m^2+n^2)^2\right).$$
We also write $D_0(t_0)$ for the integer obtained from $G(m,n)$ by removing the square factors of all primes below $10^5$. Then $E^{D(t_0)}\cong E_{g(t_0)}$ over $\Q$, and $P(t_0)$ and $Q(t_0)$ are rational points on $E^{D(t_0)}$. We note that $g(-t)=g(t)$ and $g(1/t)=g(t)/t^{14}$, so the parameters $\pm t_0^{\pm 1}$ produce the same twist, and for $t_0=\pm 1$ the formulas give $Q=\pm P$.

The polynomial $g(t)$ is the product of $-ab(t^2+1)$ and $I(t)=b^2(t^4+t^2+1)^3+a^3t^4(t^2+1)^2$. Since $b^2\approx 5.8\cdot 10^{219}$ and $a^3\approx -5.1\cdot 10^{220}$, the two terms in $I(t)$ have comparable size and opposite signs, and $I(t)$ vanishes at $t\approx \pm 1.3491$ and $t\approx \pm 0.7412$. In particular, these are the regions where $|g(t)|$, and therefore $|D(t)|$, is smallest.

\section{Root numbers of the twists}\label{sec-rootnumbers}

 We will use the following formulas of Rohrlich for local root numbers (see \cite{rohrlichvariation}).

\begin{prop}\label{prop-rootnumber}
Let $E/\Q$ be an elliptic curve and $S$ be the set of primes of bad reduction of $E$, together with $2$ and $3$. Let $D$ be a non-zero integer, and let $M$ be the part of $|D|$ that is relatively prime to all the primes in $S$. Then
$$W(E^D/\Q)=W_\infty(E^D/\Q)\cdot\prod_{p\in S}W_p(E^D/\Q)\cdot\kron{-1}{M},$$
where $W_\infty(E^D/\Q)=-1$, and $\kron{-1}{M}$ is the Jacobi symbol. Moreover, if $p\in S$ and $r_p$ is an integer in the same class as $D$ in $\Q_p^\times/(\Q_p^\times)^2$, then $W_p(E^D/\Q)=W_p(E^{r_p}/\Q)$.
\end{prop}

\begin{proof}
Let $p\not\in S$ be a prime. If $p\nmid D$, then $E^D$ has good reduction at $p$, and $W_p(E^D/\Q)=1$. If $p\mid D$, then $p\geq 5$ and $E$ has good reduction at $p$. If $v_p(D)$ is even, then $E^D$ is the twist of $E$ by a character that is unramified at $p$, so $E^D$ has good reduction at $p$, and $W_p(E^D/\Q)=1$. If $v_p(D)$ is odd, then $E^D$ has additive reduction of type $I_0^*$ at $p$, and $W_p(E^D/\Q)=\kron{-1}{p}$ by Rohrlich's formulas (\cite{rohrlichvariation}). Hence,
$$\prod_{p\not\in S}W_p(E^D/\Q)=\prod_{p\mid M,\ v_p(M) \text{ odd}}\kron{-1}{p}=\kron{-1}{M},$$
because the Jacobi symbol is multiplicative and the primes that divide $M$ to an even power contribute trivially. The last assertion holds because $E^D$ and $E^{r_p}$ are isomorphic over $\Q_p$.
\end{proof}

In practice, we choose small representatives $r_p$: if $u$ is the odd part of $D$, we let $r_2=2^{v_2(D)\bmod 2}\cdot(u\bmod 8)$, and for odd $p\in S$ we let $r_p=p^{v_p(D)\bmod 2}$ times $1$ or the least quadratic non-residue modulo $p$. Then each $W_p(E^{r_p})$ is computed directly in Magma, and the Jacobi symbol $\kron{-1}{M}$ does not require any factorization.

\begin{remark}\label{rem-closedform}
For the curve $E$ of Section \ref{sec-curve}, the local factors in Proposition \ref{prop-rootnumber} can also be written in closed form, which gives an independent check of the Magma computations. Let $\chi_D$ be the quadratic character of $\Q(\sqrt{D})$. If $E$ has multiplicative reduction at $\ell$ (that is, $\ell\in S\setminus\{5,11\}$), then by \cite{rohrlichvariation}
$$W_\ell(E^D)=\begin{cases} W_\ell(E)\,\chi_D(\ell) & \text{if } \chi_D \text{ is unramified at } \ell,\\ \chi_D(-1) & \text{if } \chi_D \text{ is ramified at } \ell.\end{cases}$$
At $\ell=5$ and $\ell=11$, where $E$ has potentially good reduction, the formulas for $\ell\geq 5$ give $W_5(E^D)=\kron{-1}{5}=1$ and $W_{11}(E^D)=\kron{-3}{11}=\kron{-1}{11}=-1$ for every $D$, because the semistability defect of $E^D$ at $5$ is $1$ or $2$, and at $11$ it is $3$ or $6$. In particular, these formulas avoid any computation of local root numbers at the wild primes $2$ and $3$ on models with very large coefficients.
\end{remark}


\begin{example}\label{ex-t3}
Let $t=3$ and $D=g(3)$. Then $D>0$ has $411$ digits, and
$$D=2^8\cdot 3^{12}\cdot 5^{15}\cdot 7^4\cdot 11^{8}\cdot 29\cdot 67\cdot 191\cdot 269\cdot 521\cdot 709\cdot 37847\cdot C,$$
where $C$ is a $362$-digit integer with no prime factors below $10^6$, which is not a perfect power, is composite (it fails a Miller--Rabin test), and is not divisible by $p$ or $q$. In particular, $D=s^2D_0$ with $s=2^4\cdot 3^6\cdot 5^7\cdot 7^2\cdot 11^4$ and $D_0=5\cdot 29\cdot 67\cdot 191\cdot 269\cdot 521\cdot 709\cdot 37847\cdot C$, which has $381$ digits, and $\Q(\sqrt{D})=\Q(\sqrt{D_0})$. The factor  $M=D/(2^8\cdot 3^{12}\cdot 5^{15}\cdot 11^8)$ of $D$ relatively prime to $S$ has $384$ digits and $M\equiv 1\bmod 4$, so $\kron{-1}{M}=1$. The local root numbers are $W_p(E^D)=-1$ for $p=11$, $31$, $41$, and $89$, and $W_p(E^D)=1$ for all other $p\in S$. This agrees with Remark \ref{rem-closedform}: $\chi_D$ is unramified at every $\ell\in S\setminus\{5\}$; the odd part of $D$ is $\equiv 5\bmod 8$, so $\chi_D(2)=-1$; and, writing $D_\ell=D/\ell^{v_\ell(D)}$, the symbol $\kron{D_\ell}{\ell}$ equals $-1$ for $\ell=3,31,37,89$ and $1$ for $\ell=13,17,23,41,43,53,p,q$. Multiplying by the values of $W_\ell(E)$ in Section \ref{sec-curve} gives the values above. Therefore, Proposition \ref{prop-rootnumber} shows that
$$W(E^D)=(-1)\cdot(-1)^4\cdot 1=-1.$$
\end{example}

\section{Independence without heights}\label{sec-independence}

The canonical heights of points on $E^D$ would require the factorization of the discriminant of $E^D$, so instead we prove independence using the $2$-descent map.

\begin{prop}\label{prop-independence}
Let $\theta$ be a root of $x^3+ax+b$, and let $K_3=\Q(\theta)$. Let $D\in\Q^\times$, let $E_D\colon y^2=x^3+aD^2x+bD^3$, and let $P,Q\in E_D(\Q)$. If $x(P)-D\theta$, $x(Q)-D\theta$, and their product are not squares in $K_3$, then $P$ and $Q$ are independent points in $E_D(\Q)$.
\end{prop}

\begin{proof}
The polynomial $x^3+ax+b$ is irreducible, so $E_D(\Q)[2]$ is trivial and every torsion point of $E_D(\Q)$ has odd order. The $2$-descent map
$$\delta\colon E_D(\Q)/2E_D(\Q)\to K_3^\times/(K_3^\times)^2,\quad (x,y)\mapsto x-D\theta,$$
is an injective homomorphism (see \cite[Ch.~X]{silverman1}), and our hypotheses say that $\delta(P)$ and $\delta(Q)$ are independent in the $\F_2$-vector space $K_3^\times/(K_3^\times)^2$. Suppose that $mP+nQ=T$, where $T$ is a torsion point and $\gcd(m,n)=1$. Since $T$ has odd order, $T\in 2E_D(\Q)$, so $mP+nQ\in 2E_D(\Q)$ and $\delta(P)^m\delta(Q)^n=1$, with $m$ or $n$ odd, which is a contradiction.
\end{proof}

To show that $\alpha=x-D\theta$ is not a square in $K_3$, it suffices to find an odd prime $\ell$ and a simple root $r$ of $x^3+ax+b$ modulo $\ell$ (which corresponds to a prime $\mathfrak{l}\mid\ell$ of degree one in $K_3$, with $\theta\equiv r\bmod\mathfrak{l}$) such that $x$ and $D$ are $\ell$-integral and $x-Dr$ is a quadratic non-residue modulo $\ell$. Indeed, $\alpha$ is then a unit at $\mathfrak{l}$ whose image in the residue field $\F_\ell$ is not a square.

\begin{example}\label{ex-t3-indep}
Let $t=3$ and $D=g(3)$. Then $x(P(3))-D\theta$ is a non-square, as witnessed by $\ell=61$ and $r=15$, and $x(Q(3))-D\theta$ and the product $(x(P(3))-D\theta)(x(Q(3))-D\theta)$ are both non-squares, as witnessed by $\ell=47$ and $r=37$. In both cases $r$ is a simple root and $\ell\nmid 90a$, so $x(P(3))$ and $x(Q(3))$, whose denominators divide $10a$ and $90a$, respectively, are $\ell$-integral. Hence, Proposition \ref{prop-independence} shows that $P(3)$ and $Q(3)$ are independent points of infinite order on $E^D$.
\end{example}

\section{The search within Mestre's family}\label{sec-search}

We have run over all $t=m/n>1$ in lowest terms with $m\leq 60$, which, by the symmetry $t\mapsto \pm t^{\pm 1}$, covers all parameters of height at most $60$, for a total of $1101$ parameters. For each $t$, we computed $\log_{10}|D_0(t)|$, where $D_0(t)$ is defined in Section \ref{sec-mestre}, and for the smallest values we computed $W(E^{D(t)}/\Q)$ via Proposition \ref{prop-rootnumber}, with the local factors at $S$ given by Remark \ref{rem-closedform} (and checked against Magma), and, when $W(E^{D(t)}/\Q)=-1$, we verified the hypotheses of Proposition \ref{prop-independence}. The ten smallest values are as follows:
\begin{center}
{\small\setlength{\tabcolsep}{3.5pt}\begin{tabular}{|c|cccccccccc|}
\hline
$t$ & $13/10$ & $53/31$ & $31/22$ & $2$ & $28/25$ & $52/29$ & $13/11$ & $11/2$ & $20/11$ & $5$ \\ \hline
$\log_{10}|D_0(t)|$ & $371.35$ & $373.14$ & $373.53$ & $374.64$ & $375.45$ & $375.46$ & $376.06$ & $376.09$ & $376.19$ & $376.20$\\\hline
$W(E^{D(t)}/\Q)$ & $1$ & $-1$ & $-1$ & $1$ & $-1$ & $-1$ & $1$ & $1$ & $-1$ & $1$\\
\hline
\end{tabular}}
\end{center}
The values of $|D_0(t)|$ range from $10^{371.35}$ to $10^{405.53}$, so the cancellation in $I(t)$ only saves a few digits: rational numbers of small height do not come much closer to the zeros of $I(t)$ than the parameters in the table. The parameter $t=3$, with $\log_{10}|D_0(3)|=380.42$, is the $28$th smallest. There are $16$ parameters with $W(E^{D(t)}/\Q)=-1$ and $|D_0(t)|<|D_0(3)|$, namely
$$t=\tfrac{53}{31},\ \tfrac{31}{22},\ \tfrac{28}{25},\ \tfrac{52}{29},\ \tfrac{20}{11},\ \tfrac{43}{29},\ \tfrac{5}{2},\ \tfrac{53}{49},\ \tfrac{19}{4},\ \tfrac{29}{17},\ \tfrac{4}{3},\ \tfrac{55}{23},\ \tfrac{11}{5},\ \tfrac{8}{5},\ \tfrac{5}{3},\ \tfrac{42}{11},$$
and for each of them the hypotheses of Proposition \ref{prop-independence} hold, so the analogue of Theorem \ref{thm-main-intro} holds with $D=D(t)$. For instance, for $t=53/31$ the integer $D_0(t)$ has $374$ digits, and the witnesses are $(\ell,r)=(61,15)$ for $P$ and $(47,37)$ for $Q$ and for the product. Among integer parameters, the smallest values are $t=2,5,12,11,17$, all with root number $1$, followed by $t=3$. We have chosen to present $t=3$ because it is the integer parameter with the smallest $|D_0(t)|$ and $W(E^{D(t)}/\Q)=-1$, and because $D=g(3)$ is then simply a value of $g$.

\section{Proof of Theorem \ref{thm-main-intro}}\label{sec-proof}


\begin{proof}
The root number of $E$ was computed in Section \ref{sec-curve}, and $W(E^D)=-1$ by Example \ref{ex-t3}. By Example \ref{ex-t3-indep}, we have $R_{E^D/\Q}\geq 2$, and therefore by Eq.~\eqref{eq-ranksum}
$$R_{E/K}=R_{E/\Q}+R_{E^D/\Q}\geq 32+2=34.$$ 
If we assume the Parity Conjecture for $E^D$, then $R_{E^D/\Q}$ is odd, because $W(E^D)=-1$, and since $R_{E^D/\Q}\geq 2$, we conclude that $R_{E^D/\Q}\geq 3$, and $R_{E/K}\geq 35$. 

Finally, if $\Sha(E^D/\Q)[p^\infty]$ is finite, then $R_{E^D/\Q}$ coincides with the $p^\infty$-Selmer corank $s_p(E^D/\Q)$, and the $p$-parity theorem (Monsky \cite{monsky} for $p=2$, and Dokchitser--Dokchitser \cite{DD10} for all $p$) shows that $(-1)^{s_p(E^D/\Q)}=W(E^D)=-1$, so the same conclusion follows.
\end{proof}

Under the isomorphism $E_D\cong E$ over $K$, which on short models is given by $(X,Y)\mapsto(X/D,Y/(D\sqrt{D}))$, the points $P(3)$ and $Q(3)$ correspond to points $P_{33},P_{34}\in E(K)$ of the form $(x,v\sqrt{D})$ with $x,v\in\Q$. The non-trivial automorphism $\sqrt{D}\mapsto -\sqrt{D}$ sends $P_{33}$ and $P_{34}$ to $-P_{33}$ and $-P_{34}$, respectively, so they span a plane in the $(-1)$-eigenspace of $E(K)\otimes\Q$, while the $32$ points $P_1,\ldots,P_{32}$ in $E(\Q)$ span a $32$-dimensional subspace of the $(+1)$-eigenspace. Hence, the points $P_1,\ldots,P_{34}$ are independent in $E(K)$. The coordinates of all these points are listed in the accompanying data files.

\begin{remark}\label{rem-selmer}
For every prime $p$, the $p$-parity theorem shows, unconditionally, that $s_p(E/\Q)$ is even and $s_p(E^D/\Q)$ is odd. Since $s_p(E/\Q)\geq R_{E/\Q}\geq 32$ and $s_p(E^D/\Q)\geq R_{E^D/\Q}\geq 2$, it follows that $s_p(E^D/\Q)\geq 3$, and therefore $s_p(E/K)=s_p(E/\Q)+s_p(E^D/\Q)\geq 35$. Here the last equality holds for every $p$, because $\Gal(K/\Q)$ acts on the $p^\infty$-Selmer group of $E/K$ tensored with $\Q_p$, and its $(\pm1)$-eigenspaces are identified with the corresponding spaces for $E/\Q$ and $E^D/\Q$.
\end{remark}

\bibliography{twist-4239}

\end{document}